% !TEX root = ../libretto.tex
% Sezione: Comunione
\momento{Comunione}
\begin{songs}{titoli}

% Fonte: "9 - Comunione/ALTISSIMO - Spoladore.pdf" (foglio accordi).
% Nell'elenco: senza tipologia.
\beginsong{Altissimo}[by={P. Spoladore}]
\beginverse
Altissimo, onnipotente, buon Signore,
tue son le lodi, la gloria, l'onore ed ogni benedizione
che a te solo e al tuo nome, Altissimo, possiamo elevare
e nessun uomo può credersi degno di poterti nominare.
\endverse
\beginverse
Laudato sii, mi Signore, con tutte le tue creature,
specialmente per frate sole così bello e radioso,
con la sua luce illumini il giorno ed illumini noi
e con grande splendore ci parla di te, Signore.
\endverse
\beginchorus
Lodate e benedite il Signore,
ringraziate e servite con grande umiltà,
lodate, lodate, benedite il Signore con grande umiltà,
ringraziate e servite con grande umiltà.
\endchorus
\beginverse
Laudato sii, mi Signore, per sora luna e le stelle,
le hai formate nel cielo così chiare, preziose e belle,
per frate vento, per l'aria, il sereno ed ogni tempo,
così la vita tu cresci e sostieni in ogni tua creatura.
\endverse
\beginverse
Laudato sii, mi Signore, per sora acqua così preziosa,
per frate fuoco giocoso e potente che c'illumina la notte,
laudato sii, mi Signore, per sora nostra madre la terra,
ci sostiene, governa e ci dona fiori, frutti ed erba.
\endverse
\beginverse
Laudato sii, mi Signore, per quelli che per il tuo amore
perdonano e sopportano in pace ogni persecuzione,
che sora morte ha trovato viventi nella tua volontà,
da te, Altissimo, un giorno saranno da te incoronati.
\endverse
\endsong

% Fonte: LibrettoAccordi.pdf (Canti Coro), J15.
% Nell'elenco: Comunione, Pentecoste; l'utente lo vuole tra i canti di Comunione.
\beginsong{Come fuoco vivo}[by={Gen Rosso, Gen Verde (CD Come fuoco vivo)}]
\beginchorus
Come fuoco vivo si accende in noi un'immensa felicità,
che mai più nessuno ci toglierà, perché tu sei ritornato.
Chi potrà tacere, da ora in poi, che sei tu in cammino con noi,
che la morte è vinta per sempre, che ci hai ridonato la vita.
\endchorus
\beginverse
Spezzi il pane davanti a noi, mentre il sole è al tramonto:
ora gli occhi ti vedono, sei tu! Resta con noi.
\endverse
\beginverse
E per sempre ti mostrerai in quel gesto d'amore:
mani che ancora spezzano pane d'eternità.
\endverse
\endsong

% Fonte: LibrettoAccordi.pdf (Canti Coro), A13.
% Nell'elenco: Comunione, Ordinario.
\beginsong{Come l'aurora verrai}[by={Gen Verde}]
\beginverse
Come l'aurora verrai, le tenebre in luce cambierai
tu per noi, Signore.
Come la pioggia cadrai, sui nostri deserti scenderai,
scorrerà l'amore.
\endverse
\beginchorus
Tutti i nostri sentieri percorrerai,
tutti i figli dispersi raccoglierai,
chiamerai da ogni terra il tuo popolo,
in eterno ti avremo con noi.
\endchorus
\beginverse
Re di giustizia sarai, le spade in aratri forgerai:
ci darai la pace.
Lupo ed agnello vedrai, insieme sui prati dove mai
tornerà la notte.
\endverse
\beginverse
Dio di salvezza tu sei e come una stella sorgerai
su di noi per sempre.
E chi non vede, vedrà, chi ha chiusi gli orecchi sentirà,
canterà di gioia.
\endverse
\endsong

% Fonte: LibrettoAccordi.pdf (Canti Coro), J36.
% Nell'elenco: Comunione, Ordinario.
\beginsong{Dietro un semplice volto}[by={M. Balduzzi (CD Giorno senza tramonto)}]
\beginchorus
Ogni volta voi lo farete a me, dietro il volto di un uomo.
Ogni volta voi scoprirete me, dietro un semplice volto.
\endchorus
\beginverse
Ogni volta che darai una goccia d'acqua,
ogni volta che sarai seme di speranza.
\endverse
\beginverse
Ma quando mai, Signore, ti abbiam visto in chi ha sete e fame?
Quando ti abbiamo visto sfigurato in chi non ha nessuno?
\endverse
\beginverse
Ogni volta che sarai una mano tesa,
ogni volta che darai forza di lottare.
\endverse
\beginverse
Ma quando mai, Signore, ti abbiam visto in chi non ha più casa?
Quando ti abbiamo visto disperato in chi chiedeva aiuto?
\endverse
\endsong

% Fonte: testo fornito dall'utente (tolti gli accordi). Non presente nell'elenco.
% Brano che dà il titolo al CD "Dove tu sei" (1982). La strofa si canta due volte, poi il finale.
\beginsong{Dove tu sei}[by={Gen Rosso (CD Dove tu sei)}]
\beginverse*
Dove tu sei torna la vita,
dove tu passi fiorisce il deserto,
dove tu guardi si rischiara il cielo
e in fondo al cuore torna il sereno.
Dove tu sei,
dove tu sei. \rep{2}
\endverse
\beginverse*
Dove tu sei
torna la vita.
\endverse
\endsong

% Fonte: LibrettoAccordi.pdf (Canti Coro), J23.
% Nell'elenco: Comunione, Ordinario.
\beginsong{E sono solo un uomo}[by={P. A. Sequeri}]
\beginverse
Io lo so, Signore, che vengo da lontano,
prima nel pensiero e poi nella tua mano,
io mi rendo conto che Tu sei la mia vita,
e non mi sembra vero di pregarti così.
\endverse
\beginverse
Padre di ogni uomo, e non ti ho visto mai,
spirito di vita, e nacqui da una donna,
figlio, mio fratello, e sono solo un uomo,
eppure io capisco che Tu sei verità.
\endverse
\beginchorus
E imparerò a guardare tutto il mondo
con gli occhi trasparenti di un bambino,
e insegnerò a chiamarti Padre nostro
ad ogni figlio che diventa uomo. \rep{2}
\endchorus
\beginverse
Io lo so Signore che Tu mi sei vicino,
Luce alla mia mente guida al mio cammino,
mano che sorregge, sguardo che perdona
e non mi sembra vero che Tu esista così.
\endverse
\beginverse
Dove nasce amore, Tu sei la Sorgente;
dove c'è una croce, Tu sei la Speranza;
dove il tempo ha fine, Tu sei Vita eterna,
e so che posso sempre contare su di Te.
\endverse
\beginchorus
E accoglierò la vita come un dono,
e avrò il coraggio di morire anch'io,
e incontro a Te verrò col mio fratello
che non si sente amato da nessuno. \rep{2}
\endchorus
\endsong

% Fonte: LibrettoAccordi.pdf (Canti Coro), J10.
% Nell'elenco: Comunione, Ordinario.
\beginsong{Fratello sole, sorella luna}[by={R. Ortolani}]
\beginverse
Dolce sentire come nel mio cuore,
ora umilmente sta nascendo amore.
Dolce capire che non son più solo
ma che son parte di una immensa vita,
che generosa risplende intorno a me;
dono di Lui del suo immenso amor.
\endverse
\beginverse
Ci ha dato il cielo e le chiare stelle,
fratello sole e sorella luna;
la madre terra, con frutti, prati e fiori,
il fuoco, il vento, l'aria e l'acqua pura,
fonte di vita per le sue creature,
dono di Lui del suo immenso amor. \rep{2}
\endverse
\endsong

% Fonte: LibrettoAccordi.pdf (Canti Coro), J31.
% Nell'elenco: Comunione, Ordinario.
\beginsong{Insieme a te}[by={F. Buttazzo}]
\beginchorus
Insieme a te, uniti a te,
la nostra vita si trasformerà.
Insieme a noi, accanto a noi,
ti sentiremo ogni giorno Gesù.
\endchorus
\beginverse
Spezzando il pane hai detto ai tuoi:
«Questo è il mio corpo donato a voi»,
prendendo il calice hai detto Gesù:
«Ecco il mio sangue versato per voi».
\endverse
\beginverse
Se celebriamo la Pasqua con te,
diventeremo discepoli tuoi,
tu hai donato la vita per noi,
perché viviamo in eterno con te.
\endverse
\beginverse
Tu hai mandato i discepoli tuoi,
in tutto il mondo a parlare di te,
a rinnovare il tuo gesto d'amore:
«Fate questo in memoria di me».
\endverse
\beginverse
E con la forza che viene da te,
cammineremo nel mondo Signor,
con questo pane che hai dato a noi,
riceveremo la vita di Dio.
\endverse
\endsong

% Fonte: LibrettoAccordi.pdf (Canti Coro), J20.
% Nell'elenco: Comunione, Ordinario.
\beginsong{L'unico maestro}[by={M. Civico}]
\beginverse
Le mie mani, con le tue, possono fare meraviglie,
possono stringere e perdonare e costruire cattedrali.
Possono dare da mangiare e far fiorire una preghiera.
\endverse
\beginchorus
Perché tu, solo tu, solo tu sei il mio maestro e insegnami
ad amare come hai fatto tu con me,
se lo vuoi, io lo grido a tutto il mondo che tu sei,
l'unico maestro sei per me.
\endchorus
\beginverse
Questi occhi, con i tuoi, potranno vedere meraviglie,
potranno piangere e luccicare, guardare oltre ogni frontiera.
Potranno amare più di ieri se sanno insieme a te sognare.
\endverse
\beginverse
Questi piedi, con i tuoi possono fare strade nuove,
possono correre e riposare, sentirsi a casa in questo mondo.
Possono mettere radici e passo a passo camminare.
\endverse
\beginverse
Tu sei il capo, noi le membra, diciamo un'unica preghiera,
Tu sei il maestro noi, testimoni della parola del vangelo.
Possiamo vivere felici in questa chiesa che rinasce.
\endverse
\endsong

% Fonte: testo fornito dall'utente (in Canti Coro c'e' solo l'immagine "1 Canti di Ingresso/Le legge della vita.jpg").
% Nell'elenco: Comunione.
\beginsong{La legge della vita}[by={Gen Rosso}]
\beginverse*
C'è una legge vera nella vita, impressa in ogni cosa:
legge che muove gli astri del cielo in un concerto d'armonia;
e in terra canta nei colori della natura, canta nella natura.
\endverse
\beginverse*
Il giorno cede il passo alla notte per amore,
la notte saluta il giorno per amore.
Dal mare sale l'acqua al cielo per amore,
e l'acqua discende dal cielo al mare per amore.
La pianta dà le foglie alla terra per amore,
la terra ridona le foglie per amore,
un seme cade in terra e muore per amore,
la vita germoglia dal solco della morte per amore!
\endverse
\beginchorus
Questa è la legge eterna, legge di Dio,
un Dio che per amore ha creato ogni cosa
ed ha nascosto amore, ed ha nascosto amore,
dietro apparenze di morte e di dolore.
\endchorus
\beginverse*
È la legge vera della vita impressa in ogni cosa,
legge che muove gli astri del cielo in un concerto d'armonia;
e in terra canta nei colori della natura, canta nella natura. \rep{2}
\endverse
\endsong

% Fonte: LibrettoAccordi.pdf (Canti Coro), A3.
% Nell'elenco: Ingresso, Ordinario (RnS); confermato dall'utente: Comunione, M. Balduzzi.
\beginsong{Nel tuo amore siamo un popolo}[by={M. Balduzzi (CD Xaipe)}]
\beginverse
Come note di una musica, pagine di un solo libro,
come il suono della pioggia, alberi di una foresta,
senza più barriere, senza più stranieri,
nel tuo amore siamo un popolo.
\endverse
\beginverse
Dove anche solo due o tre sono uniti nel tuo amore,
dove il pane della vita hai diviso insieme a noi,
senza più barriere, senza più stranieri,
nel tuo amore siamo un popolo.
\endverse
\beginverse
Quando il seme della pace spunta lieve sulla terra,
quando il soffio del tuo vento muove i giorni della storia,
senza più barriere, senza più stranieri,
nel tuo amore siamo un popolo. \rep{2}
\endverse
\endsong

% Fonte: LibrettoAccordi.pdf (Canti Coro), J8.
% Nell'elenco: Comunione, Ordinario.
\beginsong{Nel tuo silenzio}[by={Gen Verde}]
\beginverse
Nel tuo silenzio accolgo il mistero venuto a vivere dentro di me.
Sei tu che vieni, o forse è più vero che tu mi accogli in te, Gesù.
\endverse
\beginverse
Sorgente viva che nasce dal cuore, è questo dono che abita in me.
La tua presenza è un Fuoco d'amore che avvolge l'anima mia, Gesù.
\endverse
\beginverse
Ora il tuo Spirito in me dice: «Padre», non sono io a parlare, sei tu.
Nell'infinito oceano di pace Tu vivi in me, io in te, Gesù.
\endverse
\endsong

% Fonte: LibrettoAccordi.pdf (Canti Coro), J34.
% Nell'elenco: Comunione, Ordinario.
\beginsong{Nelle tue mani\soloindice{ (RnS)}}[by={RnS}]
\beginverse
Nelle tue mani affido la vita, Dio, mia salvezza sei tu.
Volgi il tuo sguardo al mio cuore, con te al sicuro sarò.
\endverse
\beginchorus
Padre del cielo, per il tuo nome vivrò,
un sacrificio con la mia lode io ti offrirò,
per tutto ciò che sempre mi dai.
Padre del cielo, per il tuo nome vivrò,
scruta il mio cuore e la mia lode io ti offrirò,
per tutto ciò che sempre mi donerai,
o Dio di ogni bontà.
\endchorus
\beginverse
Nelle tue mani è la mia vita, Dio, mia speranza sei tu.
Donami pace, o Signore, con te al sicuro vivrò.
\endverse
\endsong

% Fonte: LibrettoAccordi.pdf (Canti Coro), J5.
% Nell'elenco: Comunione (anche Ingresso).
\beginsong{Ogni mia parola}[by={Gen Verde}]
\beginverse
Come la pioggia e la neve
scendono giù dal cielo
e non vi ritornano senza irrigare
e far germogliare la terra,
così ogni mia parola non ritornerà a me
senza operare quanto desidero,
senza aver compiuto ciò per cui l'avevo mandata.
Ogni mia parola, ogni mia parola.
\endverse
\endsong

% Fonte: LibrettoAccordi.pdf (Canti Coro), J7.
% Nell'elenco: Comunione, Ordinario. CD: Dove tu sei (1982).
\beginsong{Padre mio}[by={Gen Rosso (CD Dove tu sei)}]
\beginverse
Padre mio, m'abbandono a te, di me fa' quello che ti piace;
grazie di ciò che fai per me,
spero solamente in te.
Purché si compia il tuo volere in me e in tutti i miei fratelli
niente desidero di più:
fare quello che vuoi tu.
\endverse
\beginchorus
Dammi che ti riconosca, dammi che ti possa amare sempre più.
Dammi che ti resti accanto, dammi d'essere l'amor.
\endchorus
\beginverse
Fra le tue mani depongo la mia anima, con tutto l'amore del mio cuore:
mio Dio, la dono a te,
perché ti amo immensamente.
Sì, ho bisogno di donarmi a te senza misura, affidarmi alle tue mani,
perché sei il Padre mio,
perché sei il Padre mio.
\endverse
\endsong

% Fonte: testo fornito dall'utente.
% Nell'elenco: Comunione, Ordinario. CD: Dove tu sei (1982).
\beginsong{Pane del cielo}[by={Gen Rosso (CD Dove tu sei)}]
\beginchorus
Pane del cielo sei tu, Gesù,
via d'amore: tu ci fai come te. \rep{2}
\endchorus
\beginverse
No, non è rimasta fredda la terra:
tu sei rimasto con noi
per nutrirci di te,
Pane di vita,
ed infiammare col tuo amore
tutta l'umanità.
\endverse
\beginverse
Sì, il cielo è qui su questa terra:
tu sei rimasto con noi
ma ci porti con te
nella tua casa,
dove vivremo insieme a te
tutta l'eternità.
\endverse
\beginverse
No, la morte non può farci paura:
tu sei rimasto con noi
e chi vive in te
vive per sempre.
Sei Dio per noi, sei Dio con noi,
Dio in mezzo a noi.
\endverse
\endsong

% Fonte: immagine del testo fornita dall'utente.
% Nell'elenco: Comunione, Ordinario.
\beginsong{Pane di vita}[by={Parisi}]
\beginverse
Pane di vita tu sei, Signore,
pane del cielo che nutre il cammino;
noi che veniamo alla mensa del regno
oggi con fede annunciamo il mistero.
\endverse
\beginverse
Questo è il mio corpo: prendete e mangiate;
questo è il mio sangue: prendete e bevete;
voi che accogliete la mia parola
della mia vita in eterno vivrete.
\endverse
\beginverse
Dio d'amore tu sei, o Signore,
nuova alleanza e fonte di vita;
noi ti accogliamo, o Dono del Padre,
noi ti portiamo una sete infinita.
\endverse
\beginverse
Siamo il tuo corpo che vive nel tempo,
ma non ci basta l'intero universo;
Cristo Signore, tu sei nutrimento,
placa la fame di gioia e d'eterno.
\endverse
\beginverse
Vieni tra noi e rinnova l'incontro,
pace rinnova tra il cielo e la terra;
ad ogni cosa ridona bellezza,
ad ogni uomo ridona amicizia.
\endverse
\beginverse
Figli del giorno chiamati alla luce,
nel tuo splendore viviamo, Signore;
tu sei il sole che vince la notte,
tu sei la vita che vince la morte.
\endverse
\beginverse
Popolo santo in cammino nel mondo,
tempio che accoglie la tua presenza
noi ti cantiamo la lode del cuore:
tu sei, Signore, la nostra speranza.
\endverse
\endsong

% Fonte: testo fornito dall'utente.
% Nell'elenco: Comunione, Ordinario.
\beginsong{Pane di vita nuova}[by={M. Frisina}]
\beginverse
Pane di vita nuova,
vero cibo dato agli uomini,
nutrimento che sostiene il mondo,
dono splendido di grazia.
\endverse
\beginverse
Tu sei sublime frutto
di quell'albero di vita
che Adamo non poté toccare:
ora è in Cristo a noi donato.
\endverse
\beginchorus
Pane della vita,
sangue di salvezza,
vero corpo, vera bevanda,
cibo di grazia per il mondo.
\endchorus
\beginverse
Sei l'agnello immolato
nel cui sangue è la salvezza,
memoriale della vera Pasqua
della nuova alleanza.
\endverse
\beginverse
Manna che nel deserto
nutri il popolo in cammino,
sei sostegno e forza nella prova
per la Chiesa in mezzo al mondo.
\endverse
\beginverse
Vino che ci dà gioia,
che riscalda il nostro cuore,
sei per noi il prezioso frutto
della vigna del Signore.
\endverse
\beginverse
Dalla vite ai tralci
scorre la vitale linfa
che ci dona la vita divina,
scorre il sangue dell'amore.
\endverse
\beginverse*
Amen.
\endverse
\endsong

% Testo fornito dall'utente. Parafrasi del Salmo 23 (22).
% Album: "Cerco il tuo volto" (2000).
% Struttura cantata: Rit. - Strofa - Rit. - ... (ritornello dopo ogni strofa).
% Nell'elenco: Comunione, Ordinario.
\beginsong{Perché tu sei con me}[by={Gen Verde}]
\beginchorus
Solo tu sei il mio pastore,
niente mai mi mancherà.
Solo tu sei il mio pastore, o Signore (o Signore).
\endchorus
\beginverse
Mi conduci dietro te sulle verdi alture,
ai ruscelli tranquilli lassù,
dov'è più limpida l'acqua per me,
dove mi fai riposare.
\endverse
\beginverse
Anche fra le tenebre d'un abisso oscuro
io non temo alcun male perché
tu mi sostieni, sei sempre con me,
rendi il sentiero sicuro.
\endverse
\beginverse
Siedo alla tua tavola che mi hai preparato
ed il calice è colmo per me
di quella linfa di felicità
che per amore hai versato.
\endverse
\beginverse
Sempre mi accompagnano lungo estati e inverni
la tua grazia, la tua fedeltà.
Nella tua casa io abiterò
fino alla fine dei giorni.
\endverse
\endsong

% Fonte: LibrettoAccordi.pdf (Canti Coro), J13; c'e' anche "9 - Comunione/PREGHIERA - mgc D.pdf".
% Omessa la parte finale a due voci. Nell'elenco: Comunione, Ordinario.
\beginsong{Preghiera}[by={Movimento Giovanile Costruire}]
\beginverse
Voglio ringraziarti, mio Signore, perché sei oggi ancora con me.
C'è tanta sete nel mio cuore, dammi l'acqua del tuo amore.
Sei Tu la mia certezza, sei Tu che fai sperare che
ancora il seme germoglierà, una nuova terra rivivrà,
un lavoro nuovo nascerà, tanta gente finalmente sorriderà,
un bambino forte ti pregherà, ti ringrazierà, o mio Signore.
\endverse
\beginverse
Anche quando il buio è dentro me, ed io non vedo la tua luce,
so che Tu aspetti là, e mi dai ancora la tua pace.
Fa', Signore, che non manchi mai, fa' che resti qui con noi,
facci uniti nel tuo nome, dacci ancora del tuo pane,
effondi amore su ogni cuore, dona gioia a chi è nel dolore.
Ancora il seme germoglierà, una nuova terra rivivrà,
un lavoro nuovo nascerà, tanta gente finalmente sorriderà,
ed ognuno di noi risorgerà, ed ognuno di noi per te sarà.
\endverse
\beginchorus
Preghiera, per Te sarà, preghiera. \rep{2}
\endchorus
\endsong

% Fonte: LibrettoAccordi.pdf (Canti Coro), J18.
% Nell'elenco: Comunione (autore RnS). CD: Cantiamo con gioia (1998).
\beginsong{Re di gloria}[by={F. Marranzino, A. De Luca (RnS) (CD Cantiamo con gioia)}]
\beginverse
Ho incontrato te, Gesù, e ogni cosa in me è cambiata,
tutta la mia vita ora ti appartiene.
Tutto il mio passato io lo affido a te,
Gesù, Re di gloria, mio Signor.
\endverse
\beginchorus
Tutto in te riposa, la mia mente, il mio cuore,
trovo pace in te, Signor, tu mi dai la gioia.
Voglio stare insieme a te, non lasciarti mai,
Gesù, Re di gloria, mio Signor.
\endchorus
\beginverse
Dal tuo amore chi mi separerà? Sulla croce hai dato la vita per me.
Una corona di gloria mi darai quando un giorno ti vedrò.
\endverse
\beginverse*
Io ti aspetto, mio Signor!
Io ti aspetto, mio Signor!
Io ti aspetto, mio Re!
\endverse
\endsong

% Fonte: LibrettoAccordi.pdf (Canti Coro), O22.
% Nell'elenco: Comunione, Pasqua; l'utente lo vuole tra i canti di Comunione. CD: Il mistero pasquale (2011).
\beginsong{Resto con te}[by={Gen Verde (CD Il mistero pasquale)}]
\beginverse
Seme gettato nel mondo,
Figlio donato alla terra,
il tuo silenzio custodirò.
In ciò che vive e che muore
vedo il tuo volto d'amore:
sei il mio Signore e sei il mio Dio.
\endverse
\beginchorus
Io lo so che tu sfidi la mia morte.
Io lo so che tu abiti il mio buio.
Nell'attesa del giorno che verrà
resto con te.
\endchorus
\beginverse
Nube di mandorlo in fiore
dentro gli inverni del cuore
è questo pane che tu ci dai.
Vena di cielo profondo
dentro le notti del mondo
è questo vino che tu ci dai.
\endverse
\beginverse*
Tu sei Re di stellate immensità
e sei tu il futuro che verrà,
sei l'amore che muove ogni realtà
e tu sei qui.
Resto con te.
\endverse
\endsong

% Fonte: LibrettoAccordi.pdf (Canti Coro), O20 (li' in Settimana Santa).
% Nell'elenco: Comunione, Ordinario (due righe). CD: Dove tu sei (1982).
\beginsong{Sei Dio}[by={Gen Rosso (CD Dove tu sei)}]
\beginverse
Perché avessimo la luce, ti facesti buio.
Perché avessimo la vita tu provasti la morte.
Ci basta, Signore, vederci simili a te
e offrire col tuo il nostro dolore.
\endverse
\beginchorus
Sei Dio, sei il mio Dio,
il nostro Dio
d'amore infinito. \rep{2}
\endchorus
\beginverse
Perché fossimo innocenti ti sentisti peccato.
Perché fosse nostro il cielo fosti abbandonato.
Per darci, Signore, la vera vita qui in terra,
sembrò che il Padre fosse lontano da te.
\endverse
\endsong

% Fonte: testo fornito dall'utente (tolti accordi e intermezzo). Non presente nell'elenco.
% Finale: ritornello ripetuto un tono sopra.
\beginsong{Sei il Pane della vita}[by={F. Volpintesta}]
\beginverse
A te veniamo per dirti grazie,
solo chiediamo il tuo amor,
e nasce un canto che ci innamora:
la nostra vita doniamo a te.
Sei tu soltanto il nostro bene,
vogliam restare accanto a te.
La vita è un volo che giunge a sera:
nessuno è solo, se spera in te.
\endverse
\beginchorus
Sei il Pane della vita, sei l'Amore, tu.
Sei il Giorno che non muore, sei la Vita, tu.
\endchorus
\beginverse
Ogni dolore nel tuo dolor,
ogni fatica l'hai presa tu,
nel tuo silenzio un solo canto:
ogni respiro sarà per te.
\endverse
\beginchorus
Sei il Pane della vita, sei l'Amore, tu.
Sei il Giorno che non muore, sei la Vita, tu. \rep{2}
\endchorus
\endsong

% Fonte: LibrettoAccordi.pdf (Canti Coro), J16.
% Nell'elenco: Comunione, Ordinario. CD: In cerca d'autore (1977).
\beginsong{Symbolum '77}[by={P. A. Sequeri (CD In cerca d'autore)}]
\beginverse
Tu sei la mia vita, altro io non ho; Tu sei la mia strada, la mia verità.
Nella tua parola io camminerò
finché avrò respiro, fino a quando Tu vorrai.
Non avrò paura, sai, se Tu sei con me; io Ti prego resta con me.
\endverse
\beginverse
Credo in Te, Signore, nato da Maria, Figlio eterno e santo, uomo come noi,
morto per amore, vivo in mezzo a noi,
una cosa sola con il Padre e con i tuoi;
fino a quando, io lo so, Tu ritornerai per aprirci il Regno di Dio.
\endverse
\beginverse
Tu sei la mia forza, altro io non ho; Tu sei la mia pace, la mia libertà.
Niente nella vita ci separerà.
So che la tua mano forte non mi lascerà,
so che da ogni male Tu mi libererai; e nel tuo perdono vivrò.
\endverse
\beginverse
Padre della vita, noi crediamo in Te; Figlio Salvatore, noi speriamo in Te;
Spirito d'amore, vieni in mezzo a noi.
Tu da mille strade ci raduni in unità
e per mille strade, poi, dove Tu vorrai, noi saremo il seme di Dio.
\endverse
\endsong

% Fonte: LibrettoAccordi.pdf (Canti Coro), J3; c'e' anche
% "Canti Coro/9 - Comunione/208TE, AL CENTRO DEL MIO CUORE.pdf".
% Nell'elenco: Comunione, Ordinario.
\beginsong{Te al centro del mio cuore}[by={Gen Verde}]
\beginverse
Ho bisogno di incontrarti nel mio cuore,
di trovare te, di stare insieme a te.
Unico riferimento del mio andare,
unica ragione tu, unico sostegno, tu;
al centro del mio cuore ci sei solo tu.
\endverse
\beginverse
Anche il cielo gira intorno e non ha pace,
ma c'è un punto fermo, è quella stella là:
la stella polare è fissa, ed è la sola:
la stella polare, tu, la stella sicura, tu;
al centro del mio cuore ci sei solo tu.
Tutto ruota intorno a te e in funzione di te,
e poi non importa il come, il dove, il se.
\endverse
\beginchorus
Che tu splenda sempre al centro del mio cuore,
il significato allora sarai tu.
Quello che farò sarà soltanto amore,
unico sostegno tu, la stella polare tu,
al centro del mio cuore ci sei solo tu.
\endchorus
\endsong

% Fonte: LibrettoAccordi.pdf (Canti Coro), J22.
% Nell'elenco: Comunione, Ordinario.
\beginsong{Tra voi la carità}[by={S. Zampa}]
\beginverse
Cos'è quest'aria che ci avvolge, questo mistero che ci supera,
e fa di noi un solo cuore, un'unica famiglia:
è l'amore di chi sa dare la vita, di chi non misura mai.
\endverse
\beginchorus
Ci sia tra voi la carità,
l'amore vero che mai finirà.
E tra la gente un fuoco che
il mondo intero trasformerà.
\endchorus
\beginverse
Sentirti ancora vivo tra di noi, al centro della nostra vita,
e ritrovare come nuova quell'unità più vera,
vivendo così, già sulla terra, un angolo di cielo.
\endverse
\endsong

% Fonte: LibrettoAccordi.pdf (Canti Coro), J25.
% Nell'elenco: Comunione, Ordinario.
\beginsong{Tu sei il pane di vita}[by={F. Baggio}]
\beginverse
Tu sei il pane di vita, sei la Parola di verità.
Tu sei luce del mondo, sei la certezza di eternità.
Tu, sollievo dell'anima, stella del nostro cammino.
Acqua dalla sorgente che sgorga dal cuore di Dio.
\endverse
\beginverse
Tu sei il nostro pastore che ci conduce lungo la via;
nulla c'è da temere se la tua guida mai mancherà.
Tu ci sveli nell'intimo il mistero di Dio.
Ci fai figli nel Figlio: noi, sacerdoti, profeti e re.
\endverse
\beginverse
Tu sei il nostro Maestro, sei la sapienza, unico Dio.
Sei la strada e la porta che apre il tempo all'eternità.
Tu, sostegno del debole, vincitore del male,
tu, ricchezza del povero ed eterna felicità.
\endverse
\beginverse
Tu sei pietra angolare, in te tutti siamo Chiesa di Dio.
Tu sei vite feconda e noi siamo i tralci uniti a te.
Sei l'amore più candido, sacrificio perfetto.
La bellezza che salva il mondo sei tu, Signore Gesù.
\endverse
\endsong

% Fonte: LibrettoAccordi.pdf (Canti Coro), J44.
% Nell'elenco: Comunione, Ordinario. CD: Alla tua mensa (2006).
\beginsong{Tu sei per noi}[by={P. Auricchio (CD Alla tua mensa)}]
\beginchorus
Tu sei per noi sorgente d'acqua viva,
Tu sei per noi la vita senza fine,
Tu sei per noi il pane del cielo,
Cristo Gesù Signore.
\endchorus
\beginverse
Pane del cielo, cibo di vita,
solo tu doni la gioia del cuore,
doni vigore, slancio nel dare,
spingi la Chiesa a gesti d'amore.
\endverse
\beginverse
Unico pane che nutre il mondo,
unico vino che dà speranza;
sei comunione che unisce i cuori,
sei comunione che ci fa uno.
\endverse
\beginverse
Verbo di Dio, Maestro e Signore,
guarda il tuo popolo sulle tue vie;
dona coraggio di scelte feconde,
mostra il tuo volto ad ogni uomo.
\endverse
\beginverse
Nella tua Chiesa risuona il Vangelo,
buona notizia che salva il mondo;
unico annuncio che infrange ogni buio,
vera speranza che illumina il cuore.
\endverse
\endsong

% Fonte: LibrettoAccordi.pdf (Canti Coro), J12. Omesse le indicazioni per il coro.
% Nell'elenco: Comunione, Ordinario.
\beginsong{Una città che ti assomiglia}[by={Gen Rosso}]
\beginverse
E mi fermai dove solo tu sai, dove al buio succede l'aurora, e mi calmai.
Non dubitai, niente al mondo vorrei che salire la cima del monte dove sei.
\endverse
\beginverse
Anche i passeri nel vento che volano hanno un nido ed al sicuro riposano.
Anche i piccoli del cervo hanno tane dove tornano.
\endverse
\beginchorus
E mi fermerò all'ombra dei tuoi portici;
abiterò in una grande città, che ti assomiglia e mi scalderà. \rep{2}
\endchorus
\beginverse*
E mi fermai dove solo tu sai, dove al buio succede l'aurora e mi calmai.
\endverse
\endsong

% Fonte: LibrettoAccordi.pdf (Canti Coro), J14.
% Nell'elenco: Comunione, Ordinario. CD: Verbum Panis (2000).
\beginsong{Verbum panis}[by={C. Casucci, M. Balduzzi (CD Verbum Panis)}]
\beginverse
Prima del tempo, prima ancora che la terra cominciasse a vivere
il Verbo era presso Dio.
Venne nel mondo, e per non abbandonarci in questo viaggio ci lasciò
tutto sé stesso come pane.
\endverse
\beginchorus
Verbum caro factum est, Verbum panis factum est. \rep{2}
\endchorus
\beginverse
Qui spezzi ancora il pane in mezzo a noi,
e chiunque mangerà non avrà più fame.
Qui vive la tua Chiesa intorno a Te,
dove ognuno troverà la sua vera casa.
\endverse
\beginverse
Prima del tempo, quando l'universo fu creato dall'oscurità
il Verbo era presso Dio.
Venne nel mondo, nella sua misericordia Dio ha mandato il Figlio suo,
tutto sé stesso come pane.
\endverse
\endsong

\end{songs}
